\documentclass[10pt,twocolumn]{article}
\usepackage{hubpaper}
\newcommand{\G}{\mathbb{G}}
\newcommand{\Enc}{\mathsf{Enc}}
\hubtitle{End-to-End Verifiable Online Voting\\ without a Blockchain: Homomorphic Tallying\\ and Zero-Knowledge Validity Proofs}
\begin{document}
\maketitle

\begin{abstract}
\noindent In the 2022 French legislative elections, more than $77\%$ of voters abroad voted online, but could not verify that their ballot was counted. We describe a Helios-type protocol in which (i) ballots are encrypted with additively homomorphic exponential ElGamal, so that the tally is computed on ciphertexts; (ii) each ballot carries non-interactive zero-knowledge proofs of validity; (iii) the decryption key is shared among trustees; (iv) every voter can check that their ballot is on a public bulletin board, and anyone can recompute the result. We state the security reductions, the soundness and zero-knowledge of the proofs, and a cost analysis showing that elections with $10^6$ voters are verifiable on commodity hardware, with none of the energy cost of a public blockchain.
\end{abstract}

\section{Requirements}
An election has $N$ eligible voters and a question with $K$ options; voter $i$ chooses $v_i\in\{0,1\}^K$ with $\ell\le\sum_jv_{ij}\le u$. We want: \emph{correctness} (the published tally is $T_j=\sum_iv_{ij}$); \emph{ballot privacy} (no coalition below a threshold learns more about $v_i$ than what $T$ reveals); \emph{individual verifiability} (voter $i$ can check that their ballot was recorded); \emph{universal verifiability} (anyone can check that $T$ is the tally of the recorded valid ballots); \emph{eligibility} (one ballot per registered voter). A central server must not be trusted for any of these. We follow Helios~\cite{helios2008}, itself built on~\cite{cgs1997}.

\section{Homomorphic encryption}
Let $\G=\langle g\rangle$ be a group of prime order $q$ in which the decisional Diffie--Hellman (DDH) problem is hard~\cite{boneh1998} (a subgroup of $\Z_p^*$, $p=2q+1$, or an elliptic curve). The secret key is $x\in\Z_q$, the public key $y=g^x$. Exponential ElGamal~\cite{elgamal1985} encrypts $m\in\Z_q$ as
\begin{equation}
\Enc(m;r)=(\alpha,\beta)=(g^r,\;g^my^r),\qquad r\xleftarrow{\$}\Z_q .
\end{equation}
\begin{proposition}[homomorphism]
$\Enc(m_1;r_1)\cdot\Enc(m_2;r_2)=\Enc(m_1+m_2;r_1+r_2)$ componentwise.
\end{proposition}
\noindent The proof is immediate. Hence $\prod_i\Enc(v_{ij};r_{ij})=\Enc(T_j;R_j)$, $R_j=\sum_ir_{ij}$: the tally is computed without decrypting any ballot. Decryption gives $g^{T_j}=\beta/\alpha^x$; since $0\le T_j\le N$, $T_j$ is recovered by baby-step giant-step~\cite{shanks1971} in $O(\sqrt N)$ group operations, i.e.\ about $10^3$ for $N=10^6$.

\textbf{Privacy.} Under DDH, ElGamal is IND-CPA: $(g^r,g^my^r)$ is indistinguishable from a pair of uniform elements. IND-CPA alone is not enough, since an adversary could copy and re-randomise someone else's ballot; attaching the proofs below makes the ballot non-malleable, and this is what yields ballot privacy in the random oracle model~\cite{bernhard2012}.

\section{Zero-knowledge validity proofs}
A malicious voter could encrypt $v_{ij}=1000$. Each ballot therefore proves, without revealing $v_i$, that every $v_{ij}\in\{0,1\}$ and that $\sum_jv_{ij}\in[\ell,u]$.

\textbf{Chaum--Pedersen.} To prove that $(\alpha,\beta)$ encrypts a given $m$, i.e.\ $\log_g\alpha=\log_y(\beta g^{-m})=r$: the prover sends $(a,b)=(g^w,y^w)$, receives a challenge $c\in\Z_q$, answers $z=w+cr$; the verifier checks
\begin{equation}
g^z=a\,\alpha^c,\qquad y^z=b\,(\beta g^{-m})^c .
\end{equation}
\begin{proposition}
This $\Sigma$-protocol~\cite{chaum1992} is complete, special sound and honest-verifier zero-knowledge.
\end{proposition}
\begin{proof}
Completeness: $g^{w+cr}=g^w(g^r)^c$. Special soundness: two accepting transcripts $(a,b,c,z)$, $(a,b,c',z')$, $c\neq c'$, give $g^{z-z'}=\alpha^{c-c'}$, hence the witness $r=(z-z')(c-c')^{-1}\bmod q$ satisfies both relations. HVZK: choose $c,z$ uniformly and set $a=g^z\alpha^{-c}$, $b=y^z(\beta g^{-m})^{-c}$; this has exactly the distribution of real transcripts.
\end{proof}
\textbf{Disjunction.} To prove ``$m=0$ or $m=1$'' without revealing which, the prover simulates the false branch with a self-chosen challenge $c_{1-m}$, runs the true branch honestly with $c_m=c-c_{1-m}$, and the verifier checks both branches and $c_0+c_1=c$~\cite{cds1994}. The same technique applied to $\prod_j(\alpha_j,\beta_j)=\Enc(\sum_jv_{ij})$ proves $\sum_jv_{ij}\in\{\ell,\dots,u\}$. If a ballot is invalid, special soundness implies that a cheating prover succeeds with probability at most $1/q\approx2^{-256}$.

\textbf{Non-interactivity.} With Fiat--Shamir~\cite{fiatshamir1986}, $c=\mathcal H(\text{statement},a,b)$ for a hash function $\mathcal H$ modelled as a random oracle. The statement ($g,y,\alpha,\beta$, election identifier) \emph{must} be hashed: the ``weak'' variant that hashes only $(a,b)$ is unsound and broke the ballot privacy of early Helios versions~\cite{bernhard2012}.

\section{Distributed decryption}
The key is shared among $n$ trustees: $x=\sum_kx_k$, $y=\prod_kg^{x_k}$ (or a Shamir $t$-of-$n$ sharing~\cite{shamir1979,pedersen1991}, so that $n-t$ trustees may be absent). Each trustee publishes $d_{jk}=\alpha_j^{x_k}$ with a Chaum--Pedersen proof that $\log_g y_k=\log_{\alpha_j}d_{jk}$, and $g^{T_j}=\beta_j/\prod_kd_{jk}$. A trustee who lies is detected with probability $1-1/q$; privacy holds unless $t$ trustees collude, and no single server, hacked or dishonest, can decrypt an individual ballot.

\section{Verification}
\textbf{Individual.} The voter keeps $h_i=\mathcal H(\text{ballot}_i)$ and checks that $h_i$ appears on the public append-only bulletin board; with collision-resistant $\mathcal H$, no other ballot can match. \emph{Cast-as-intended} is checked with the Benaloh challenge~\cite{benaloh2006}: before casting, the voter may ask the device to reveal $r$ and check the encryption, then re-encrypt. A device that cheats on a fraction $f$ of ballots, each audited with probability $p$, escapes detection on $k$ ballots with probability $(1-pf)^k$.

\textbf{Universal.} Anyone can recompute $\prod_i(\alpha_{ij},\beta_{ij})$ from the board, check all validity and decryption proofs, and hence the tally: the result is correct as soon as all proofs verify, whoever ran the servers.

\textbf{Eligibility.} The board publishes a commitment per registered voter (a keyed hash of the identifier, not a plain hash of an e-mail, which can be reversed by dictionary search) and the number of ballots; any citizen can check that their commitment is listed and that $\#\text{ballots}\le\#\text{registered}$, which bounds ballot stuffing.

\section{Cost and comparison}
On an elliptic curve, a ballot with $K$ options costs $2K$ exponentiations to encrypt and $O(K)$ more for the proofs; verifying a ballot also costs $O(K)$. For $N=10^6$ and $K=3$, universal verification is of the order of $10^7$ exponentiations, i.e.\ minutes on a laptop, and it is embarrassingly parallel. A public blockchain (e.g.\ V8TE on Ethereum) provides an append-only board at a much higher monetary and energy cost; here the board only needs to be replicated and signed, since its integrity is checked by the hashes and proofs rather than by consensus.

\textbf{Limits.} The protocol is not receipt-free: a voter can prove their vote using $r$, which enables vote buying; coercion resistance requires credential-based schemes~\cite{jcj2005}. Privacy is computational (DDH), not everlasting.

\section{Status}
A prototype based on the Helios code base covers key generation, encryption, homomorphic tally, disjunctive proofs, trustee decryption and verification; next steps are threshold key generation, the Benaloh audit and large-scale benchmarks up to $10^6$ voters.

\footnotesize
\begin{thebibliography}{99}\setlength{\itemsep}{0pt}
\bibitem{benaloh2006} J.~Benaloh. Simple verifiable elections. \emph{USENIX EVT}, 2006.
\bibitem{bernhard2012} D.~Bernhard, O.~Pereira, B.~Warinschi. How not to prove yourself: pitfalls of the Fiat--Shamir heuristic and applications to Helios. \emph{ASIACRYPT}, 2012.
\bibitem{boneh1998} D.~Boneh. The decision Diffie--Hellman problem. \emph{ANTS-III}, LNCS 1423, 1998.
\bibitem{chaum1992} D.~Chaum, T.~Pedersen. Wallet databases with observers. \emph{CRYPTO}, 1992.
\bibitem{cds1994} R.~Cramer, I.~Damg{\aa}rd, B.~Schoenmakers. Proofs of partial knowledge and simplified design of witness hiding protocols. \emph{CRYPTO}, 1994.
\bibitem{cgs1997} R.~Cramer, R.~Gennaro, B.~Schoenmakers. A secure and optimally efficient multi-authority election scheme. \emph{EUROCRYPT}, 1997.
\bibitem{elgamal1985} T.~ElGamal. A public key cryptosystem and a signature scheme based on discrete logarithms. \emph{IEEE Trans.\ Inf.\ Theory} 31(4):469--472, 1985.
\bibitem{fiatshamir1986} A.~Fiat, A.~Shamir. How to prove yourself: practical solutions to identification and signature problems. \emph{CRYPTO}, 1986.
\bibitem{helios2008} B.~Adida. Helios: web-based open-audit voting. \emph{USENIX Security}, 2008.
\bibitem{jcj2005} A.~Juels, D.~Catalano, M.~Jakobsson. Coercion-resistant electronic elections. \emph{ACM WPES}, 2005.
\bibitem{pedersen1991} T.~Pedersen. A threshold cryptosystem without a trusted party. \emph{EUROCRYPT}, 1991.
\bibitem{shamir1979} A.~Shamir. How to share a secret. \emph{Comm.\ ACM} 22(11):612--613, 1979.
\bibitem{shanks1971} D.~Shanks. Class number, a theory of factorization and genera. \emph{Proc.\ Symp.\ Pure Math.} 20:415--440, 1971.
\end{thebibliography}
\end{document}
